\SourceLeaf{255}%
\Needspace{7\baselineskip}
\begin{center}
\textsc{Gauss} to \textsc{Gerling}. Göttingen, 23 June 1846.
\end{center}

\mbox{.\ .\ .\ .\ .\ .\ .} In the last analysis the distinction between screws, like the distinction between right and left, depends on the fact that

\begin{quote}
three straight lines issuing from a point $A$ and lying in no single plane, namely $AB$, $AC$, $AD$, can, when paired with three others $ab$, $ac$, $ad$ in the order in which they are listed,
\end{quote}

\begin{center}
have the same or the opposite orientation.
\end{center}

\Needspace{14\baselineskip}
In the first case, namely, the pyramid $ABCD$ can pass into the pyramid $abcd$ by gradual quantitative changes alone, so that $A$ passes into $a$, $B$ into $b$, and so on. In the second case such a transition is possible only if the changing pyramid \textls{vanishes} at some point in between ($ABCD$ in one plane). The following have like orientations

\[
\begin{array}{ccc}
AB,&AC,&AD\\
AC,&AD,&AB\\
AD,&AB,&AC,
\end{array}
\]

whereas $AB$, $AD$, $AC$ and so on have the opposite orientation.

This \textls{difference} between two systems of three straight lines each can \textls{not} be reduced to concepts. Only by referring to \textls{actually existing} spatial things can it be \textls{shown}. Two minds can communicate about it only by attaching their intuitions to one and the same system found in the real world. If we choose for the 1\textsuperscript{st} system the [order:]

\Needspace{7\baselineskip}
\begin{center}
South \qquad West \qquad Zenith,
\end{center}

then for every other system of three it is settled whether it has the same or the opposite orientation as this one. The following agree with it:

\Needspace{7\baselineskip}
\begin{center}
\begin{tabular}{cccc}
Front & Right & Up & \\
Left & Front & Up & \\
Front & Left & Down & \\
Right & Front & Down & \\
Right & Back & Up & and so on.
\end{tabular}
\end{center}
